\documentclass[conference]{IEEEtran}
\IEEEoverridecommandlockouts
\usepackage{cite}
\usepackage{amsmath,amssymb,amsfonts}
\usepackage{algorithmic}
\usepackage{graphicx}
\usepackage{textcomp}
\usepackage{xcolor}
\usepackage{array}
\usepackage{multirow}
\def\BibTeX{{\rm B\kern-.05em{\sc i\kern-.025em b}\kern-.08em
    T\kern-.1667em\lower.7ex\hbox{E}\kern-.125emX}}
\begin{document}

\title{Smart Wardrobe: An Event-Aware Outfit Recommendation System Using ResNet50 and CLIP Embeddings}

\author{\IEEEauthorblockN{Sameer Saxena}
\IEEEauthorblockA{\textit{NUID: 002523266} \\
\textit{Group 10}\\
Khoury College of Computer Sciences \\
Northeastern University\\
saxena.same@northeastern.edu}
}

\maketitle

\begin{abstract}
This paper presents Smart Wardrobe, a locally deployed outfit recommendation system that combines visual compatibility scoring with event-aware and gender-aware classification. The system accepts user-uploaded clothing images through a Streamlit web interface, extracts 2560-dimensional embeddings by combining a pretrained ResNet50 backbone (2048-dim) with a CLIP ViT-B/32 encoder (512-dim), and stores item representations alongside gender metadata in a SQLite database. Outfit recommendations are generated by enumerating all valid top-bottom-shoes combinations and scoring each using two signals, namely a visual compatibility score from pairwise cosine similarity and an event suitability score from a trained multilayer perceptron. The event classifier takes a 2563-dimensional input and is trained on the wardrobe\_v2 dataset, a curated subset of 5,994 clothing items across four usage classes. With a 70/15/15 stratified split and inverse-frequency class weighting, the classifier achieves 88\% test accuracy and a macro AUC-ROC of 0.9631. The system is fully local with no external API dependencies, and this paper also surveys prior work in fashion recommendation, visual embeddings, and outfit compatibility to contextualise the design choices made.
\end{abstract}

\begin{IEEEkeywords}
outfit recommendation, ResNet50, CLIP, visual embeddings, event classification, gender-aware recommendation, fashion compatibility, cosine similarity, Streamlit, SQLite, deep learning
\end{IEEEkeywords}

% ─────────────────────────────────────────────────────────────
\section{Introduction}
% ─────────────────────────────────────────────────────────────

\subsection{Background}

Choosing an appropriate outfit for a specific occasion is something most people do without much thought, but it actually involves a surprising number of judgements. You have to think about what you own, whether the pieces go together visually, and whether the combination is appropriate for where you are going. As wardrobes grow larger and life gets busier, making these decisions consistently well becomes harder. This is the kind of problem that machine learning is well suited to help with.

The broader field of fashion AI has grown significantly over the last decade. Researchers have tackled problems like outfit compatibility, personalised clothing retrieval, clothing segmentation, and virtual try-on. Most of these systems are built for large e-commerce platforms and focus on helping users discover new items to buy rather than making better use of what they already own. The gap between what the research community has built and what an individual user actually needs is still quite wide.

\subsection{Objective}

The goal of this work is twofold. First, it surveys the existing literature on ML-based fashion recommendation, visual embeddings for clothing, and outfit compatibility learning, with the aim of understanding what has been tried, what has worked, and what gaps remain. Second, it presents the Smart Wardrobe system as a practical implementation that directly addresses one of those gaps, namely composing outfits from a personal wardrobe for a specific social occasion while respecting gender context.

\subsection{Scope}

The literature survey covers papers published between 2015 and 2024, selected from IEEE Xplore, ACM Digital Library, Google Scholar, and arXiv. Papers were included if they were directly relevant to at least one of the following themes: visual embedding extraction for clothing, outfit compatibility modelling, personalised fashion retrieval, or fashion classification datasets. Papers focused solely on clothing segmentation or virtual try-on were excluded since they do not directly inform the recommendation pipeline built here.

\subsection{Problem Statement}

Given a personal wardrobe of clothing items and a target event type such as office or wedding, the system must identify the top-N outfit combinations where each outfit has one top, one bottom, and one pair of shoes, that are both visually compatible and appropriate for the event.

There are two interacting sub-problems here. The first is visual compatibility, meaning do the selected items actually look good together in terms of colour, texture, and style. The second is event suitability, meaning are the items appropriate for the target occasion. An additional practical constraint is gender coherence, since outfit combinations should not mix items intended for different genders unless those items are marked as Unisex.

\subsection{Contributions}

The main contributions of this work are the following. A fully local recommendation pipeline requiring no external APIs. A hybrid 2560-dimensional embedding combining ResNet50 and CLIP that outperforms either backbone alone. A gender-aware filtering gate that ensures outfit coherence across Men, Women, and Unisex items. An automatic embedding migration system that re-embeds stored wardrobe items if the model backbone changes, avoiding manual re-uploads. A two-signal scoring framework combining unsupervised cosine similarity with a supervised MLP event classifier. And finally a working Streamlit application that brings all of these components together in a usable interface.

% ─────────────────────────────────────────────────────────────
\section{Methodology}
% ─────────────────────────────────────────────────────────────

\subsection{Search Strategy}

The literature for this survey was gathered from four main sources. IEEE Xplore was the primary database for conference papers in computer vision and machine learning. The ACM Digital Library was used for multimedia and retrieval papers. Google Scholar was used for broader coverage, particularly for preprints and papers that appeared at workshops. arXiv was consulted for more recent work that had not yet been formally published at the time of writing.

Searches were conducted between January and March 2026. The following keywords and phrases were used in various combinations: fashion recommendation, outfit compatibility, clothing embeddings, visual fashion retrieval, ResNet clothing features, CLIP fashion, personalised outfit recommendation, fashion product images, deep learning wardrobe, and occasion-aware clothing.

\subsection{Selection Criteria}

Papers were included if they met at least one of the following criteria. They proposed or evaluated a method for learning visual representations of clothing items. They addressed the problem of outfit compatibility or recommendation from a machine learning perspective. They introduced or analysed a publicly available fashion dataset relevant to compatibility or usage classification. Papers were excluded if they focused exclusively on clothing segmentation, 3D garment simulation, or virtual try-on without any retrieval or recommendation component. Papers with fewer than 20 citations that did not introduce a new dataset or method directly relevant to the current system were also excluded.

In total, eight papers were selected for detailed review. These are cited as \cite{b1} through \cite{b8} in this report.

% ─────────────────────────────────────────────────────────────
\section{Literature Review}
% ─────────────────────────────────────────────────────────────

\subsection{Theme 1: Visual Embedding Backbones for Clothing}

The foundation of most modern fashion recommendation systems is a pretrained convolutional neural network used as a fixed feature extractor. He et al. \cite{b1} introduced ResNet in 2016, demonstrating that residual skip connections could train networks of over 100 layers without the vanishing gradient problem that had limited deeper architectures. The ResNet50 variant in particular became a standard backbone across vision tasks. Its global average pooling output produces a 2048-dimensional vector that captures both low-level features like colour and texture and higher-level structural patterns. Because it is pretrained on ImageNet, these features generalise well to domains like fashion even without any fine-tuning on clothing data.

Several years later, Radford et al. \cite{b8} introduced CLIP, which took a fundamentally different training approach. Rather than training on labelled image classification benchmarks, CLIP was trained on 400 million image-text pairs using contrastive learning, where matching images and text descriptions were pulled together in embedding space while non-matching pairs were pushed apart. The result is an image encoder that understands semantic concepts rather than just visual patterns. When applied to clothing, CLIP embeddings reflect things like ``this looks formal'' or ``this is sportswear'' in a way that a ResNet trained on ImageNet categories like ``dog'' and ``car'' does not naturally capture.

These two papers represent two complementary philosophies for visual representation. ResNet excels at low-level visual features, while CLIP captures high-level semantic meaning. Using only one of them means leaving information on the table, which is the motivation for the hybrid embedding approach taken in this project.

\subsection{Theme 2: Outfit Compatibility Learning}

The most influential early work on learning outfit compatibility came from Han et al. \cite{b2}, who proposed using bidirectional LSTMs to model compatibility across sequentially ordered clothing items. Their key insight was that outfit compatibility is not just about pairwise similarity but about how items relate to each other as a set. They also introduced the Polyvore dataset, which contained 21,889 human-curated outfits with binary compatibility labels. This dataset became the standard benchmark for outfit compatibility research for several years.

Vasileva et al. \cite{b3} built on this by arguing that compatibility is not a single global concept but depends on the types of items being compared. A shirt and jeans compatibility is judged differently from a shoe and bag compatibility. Their solution was to learn type-conditioned embedding subspaces, where a separate projection matrix is learned for each pair of item types. This type-aware approach significantly improved compatibility prediction on Polyvore compared to global embedding methods.

Tangseng et al. \cite{b4} addressed a slightly different problem that is actually the closest to what Smart Wardrobe does. They focused on recommending outfits specifically from a personal closet rather than from a large catalogue, training a neural outfit grader on Polyvore outfit data and then applying it to enumerate and score combinations from a user's own items. The key difference from Han et al. and Vasileva et al. is the inference setting: instead of retrieval from a large item pool, all possible combinations are generated from a fixed personal wardrobe and ranked.

\subsection{Theme 3: Personalised Fashion Retrieval}

He and McAuley \cite{b5} approached fashion recommendation from a collaborative filtering perspective rather than a visual compatibility one. Their VBPR model extended Bayesian Personalised Ranking by incorporating visual product features extracted from images, allowing the model to make recommendations for items the user had never interacted with before. This was an important contribution because it showed how visual features could improve recommendation quality even in standard retrieval settings, not just in compatibility-focused systems. The limitation is that VBPR requires user interaction history to personalise recommendations, which is something the current Smart Wardrobe system does not yet support.

\subsection{Theme 4: Fashion Datasets}

Aggarwal \cite{b6} published the Fashion Product Images dataset sourced from Myntra, which contains over 44,000 clothing items with detailed metadata including category, gender, color, season, and a usage label covering classes like Casual, Formal, Sports, and Ethnic. Unlike Polyvore, this dataset does not have outfit-level compatibility labels. What it does have is item-level occasion annotations, which is exactly what is needed to train an event classifier. The dataset is also publicly available on Kaggle without access restrictions, which was a practical advantage over Polyvore.

\subsection{Comparative Analysis}

Comparing these papers reveals a clear progression. Han et al. and Vasileva et al. both require curated outfit-level labels for supervised training, which limits applicability to scenarios where such labels are available. Tangseng et al. applied this to the personal closet setting but still relied on Polyvore outfit data. None of these systems account for occasion or gender at all. They also require the Polyvore dataset, which has become increasingly difficult to access.

VBPR and similar collaborative filtering approaches require user history and are therefore not applicable to a cold-start wardrobe system where no prior feedback exists.

The Fashion Product Images dataset fills a different role. It does not have compatibility labels, but it does have the occasion metadata needed to supervise an event classifier. The tradeoff is that the supervision is at the item level rather than the outfit level, which introduces approximation error when aggregating item-level predictions to outfit-level scores.

Table~\ref{tab:comparison} summarises how the key related works compare across the dimensions most relevant to Smart Wardrobe.

\begin{table}[htbp]
\caption{Comparison of Related Works}
\begin{center}
\renewcommand{\arraystretch}{1.2}
\begin{tabular}{|l|c|c|c|c|}
\hline
\textbf{Work} & \textbf{Personal Closet} & \textbf{Event Labels} & \textbf{Gender} & \textbf{Open Data} \\
\hline
Han et al. \cite{b2}      & No  & No  & No  & Limited \\
Vasileva et al. \cite{b3} & No  & No  & No  & Limited \\
Tangseng et al. \cite{b4} & Yes & No  & No  & Limited \\
He \& McAuley \cite{b5}   & No  & No  & No  & Yes \\
FPI Dataset \cite{b6}     & N/A & Yes & Yes & Yes \\
Smart Wardrobe (ours)     & Yes & Yes & Yes & Yes \\
\hline
\end{tabular}
\label{tab:comparison}
\end{center}
\end{table}

% ─────────────────────────────────────────────────────────────
\section{Critical Analysis}
% ─────────────────────────────────────────────────────────────

\subsection{Evaluation of Research Quality}

The papers reviewed in this survey are generally of high quality. Han et al. \cite{b2} and Vasileva et al. \cite{b3} were published at top-tier venues (ACM Multimedia and ECCV respectively) and have accumulated thousands of citations. Their experimental setups are rigorous, with clear train/test splits on Polyvore and well-defined evaluation metrics including AUC for compatibility prediction. Tangseng et al. \cite{b4} is a workshop paper and is therefore less heavily cited, but the problem formulation is well-suited to real-world application and the experimental results on Polyvore are credible.

He et al. \cite{b1} on ResNets is one of the most cited papers in the history of computer vision, and its validity has been replicated extensively across domains. Radford et al. \cite{b8} on CLIP is similarly well-established, with the CLIP embeddings having been validated on dozens of downstream tasks including fashion understanding. He and McAuley \cite{b5} on VBPR is a solid recommender systems paper with reproducible results on standard e-commerce datasets.

The Fashion Product Images dataset \cite{b6} is less formally validated as a research contribution since it is a Kaggle dataset rather than a peer-reviewed publication. The usage labels appear to be sourced directly from Myntra product metadata, which means they reflect retailer categorisation rather than ground-truth human judgement. This introduces some label noise, particularly for items that could reasonably belong to multiple categories.

\subsection{Identification of Gaps}

Several important gaps emerge from reviewing this body of work.

The most significant gap is the complete absence of occasion or event labels in any of the outfit compatibility datasets. Polyvore outfits are labelled as compatible or not, but there is no information about what occasion they are intended for. This means none of the compatibility models from Han et al., Vasileva et al., or Tangseng et al. can differentiate between a casual outfit and a formal one at prediction time.

A second gap is the lack of gender consideration in any of the compatibility systems. Polyvore-based models treat outfit items as interchangeable regardless of whether they are men's or women's clothing. In practice this is a significant limitation because most wardrobes are strongly gendered.

Third, all of the Polyvore-based methods require access to the Polyvore dataset, which is no longer openly distributed. This is a reproducibility concern and a practical barrier to building on this line of work.

Fourth, none of the reviewed systems are designed for local deployment or privacy preservation. They assume access to server-side infrastructure, user data, and in some cases external APIs.

\subsection{Implications}

These gaps directly informed the design of Smart Wardrobe. The absence of occasion labels in compatibility datasets meant that the system could not be built as a standard compatibility classifier trained on Polyvore. Instead, an event classifier trained on the Fashion Product Images usage labels was used as a proxy for occasion suitability, combined with unsupervised cosine similarity for compatibility. This is an approximation, but it is the best available approach given the data that exists.

The absence of gender consideration in prior work motivated the explicit gender tagging system and the gender coherence filter in Smart Wardrobe. The privacy and deployment concerns motivated the local-only architecture with no external API dependencies.

The reliance on Polyvore in prior work also suggests that the field would benefit significantly from new open datasets that include both outfit-level compatibility labels and occasion annotations. The wardrobe\_v2 dataset used here is a step in this direction but does not yet include outfit-level labels.

\subsection{Limitations of This Survey}

This survey is limited in several ways. The paper selection, while systematic, is not exhaustive. Fashion AI is a large and fast-moving field and many relevant papers from 2022 to 2024, particularly those using transformer-based vision models or large multimodal models for fashion understanding, were outside the scope of this review. The focus on English-language papers from Western venues may also underrepresent work from Asian research communities where fashion AI has seen significant investment. Additionally, the comparative table in Section III is based on the papers' stated problem settings and may not fully reflect what each system could be adapted to do with additional engineering.

% ─────────────────────────────────────────────────────────────
\section{Dataset}
% ─────────────────────────────────────────────────────────────

\subsection{wardrobe\_v2 Dataset}

The event classifier is trained on the wardrobe\_v2 dataset, a curated subset of the Fashion Product Images dataset \cite{b6} from Myntra. It contains 6,000 balanced product listings with 2,000 tops, 2,000 bottoms, and 2,000 shoes. Each item has the following metadata fields: \texttt{id}, \texttt{file\_name}, \texttt{category}, \texttt{gender}, \texttt{subCategory}, \texttt{articleType}, \texttt{baseColour}, and \texttt{usage}.

The original dataset has eight usage categories. Four low-sample classes were merged before training. Smart Casual and Party were absorbed into Formal and Casual respectively, and Home and Travel were both merged into Casual. After merging, 5,994 valid items remain across four classes. Casual accounts for 4,095 items (68.3\%), Sports for 974 (16.2\%), Formal for 498 (8.3\%), and Ethnic for 427 (7.1\%). Gender distribution is Men at 60.3\%, Women at 37.6\%, and Unisex at 2.1\%.

The Casual class dominates at 68.3\% of samples, which creates a class imbalance problem. This is addressed through inverse-frequency class weighting in the loss function and \texttt{WeightedRandomSampler}, which up-samples minority classes during each training epoch.

\subsection{Data Split}

The dataset uses a stratified 70/15/15 train/val/test split giving 4,195 training, 899 validation, and 900 test samples. Stratification preserves class proportions across all three splits.

% ─────────────────────────────────────────────────────────────
\section{System Architecture}
% ─────────────────────────────────────────────────────────────

\subsection{Overview}

Smart Wardrobe is built from six modules. These are the hybrid visual embedding extractor, the SQLite wardrobe store with an embedding migration gate, the cosine similarity compatibility scorer, the MLP event classifier, the gender-aware filtering and ranking component, and the Streamlit web interface. Items are embedded once at upload time and stored. At recommendation time, gender filtering is applied, then all valid combinations are scored in a single batched pass.

\subsection{Hybrid Visual Embedding Extraction}

Each clothing image is processed through two backbones whose outputs are concatenated into a 2560-dimensional vector.

The first backbone is a pretrained ResNet50 \cite{b1} with the final classification layer removed, run in evaluation mode with no gradient computation. Input images are resized to $224\times224$ and normalised using ImageNet statistics ($\mu=[0.485, 0.456, 0.406]$, $\sigma=[0.229, 0.224, 0.225]$). The global average pooling output gives a 2048-dim embedding capturing colour, texture, and shape.

The second backbone is the CLIP ViT-B/32 image encoder \cite{b8}, loaded via \texttt{clip.load("ViT-B/32")} and run in evaluation mode. CLIP's own preprocessing is applied and the image encoding is L2-normalised before concatenation, giving a 512-dim semantic embedding.

The two are combined as

\begin{equation}
\mathbf{e} = [\mathbf{e}_\text{ResNet} \;\|\; \mathbf{e}_\text{CLIP}] \;\in\; \mathbb{R}^{2560}
\label{eq:hybrid}
\end{equation}

This fusion preserves both the low-level visual sensitivity of ResNet50 and the semantic style awareness of CLIP.

\subsection{Wardrobe Storage and Embedding Migration}

Items are stored in a SQLite database with the schema shown below.

\begin{verbatim}
CREATE TABLE items (
  id        INTEGER PRIMARY KEY,
  name      TEXT NOT NULL,
  category  TEXT NOT NULL,
  gender    TEXT NOT NULL DEFAULT 'Unisex',
  embedding BLOB NOT NULL,
  thumbnail BLOB,
  added_at  TIMESTAMP DEFAULT CURRENT_TIMESTAMP
);
\end{verbatim}

Embeddings are serialised as 32-bit float byte arrays. A $100\times100$ JPEG thumbnail is stored as a BLOB for UI display. Original images are discarded after embedding extraction to preserve privacy.

On every application startup, the system checks all stored embeddings against the required dimension of 2560. Items with mismatched dimensions are automatically re-embedded using their stored thumbnail via the current extractor. This migration gate ensures that backbone upgrades do not require users to re-upload their wardrobe.

\subsection{Visual Compatibility Scoring}

Compatibility for a three-item outfit with top $t$, bottom $b$, and shoes $s$ is computed from pairwise cosine similarities as

\begin{equation}
\text{avg\_sim} = \frac{1}{3}\bigl[\cos(t,b) + \cos(t,s) + \cos(b,s)\bigr]
\label{eq:avg_sim}
\end{equation}

\begin{equation}
\text{compat} = \frac{\text{avg\_sim} + 1}{2} \;\in\; [0,1]
\label{eq:compat}
\end{equation}

Equation~\eqref{eq:compat} maps the cosine range from $[-1,1]$ to $[0,1]$. For batch scoring, row-wise cosine similarity is applied across stacked embedding matrices

\begin{equation}
\cos_{\text{row}}(A, B)_i = \frac{A_i \cdot B_i}{\|A_i\| \cdot \|B_i\| + \epsilon}
\label{eq:rowcos}
\end{equation}

where $\epsilon = 10^{-8}$ prevents division by zero.

\subsection{Event Suitability Classification}

Gender is encoded as a 3-dimensional one-hot vector over Men, Women, and Unisex. This is concatenated with the 2560-dim embedding to give the 2563-dim classifier input

\begin{equation}
\mathbf{x} = [\mathbf{e} \;\|\; \text{onehot}(g)] \;\in\; \mathbb{R}^{2563}
\label{eq:input}
\end{equation}

The event classifier is a feedforward network mapping this to a distribution over four usage classes

\begin{equation}
f_\theta: \mathbb{R}^{2563} \rightarrow \mathbb{R}^{4}
\label{eq:mlp}
\end{equation}

The network is structured as follows. A linear layer from 2563 to 512 is followed by batch normalisation, ReLU activation, and dropout at rate 0.3. Then a second linear layer from 512 to 128 with batch normalisation, ReLU, and dropout at rate 0.2. Finally a linear layer from 128 to 4 with softmax output. Training uses the Adam optimiser \cite{b7} with learning rate $10^{-3}$ and weight decay $10^{-4}$. The learning rate is halved when validation loss plateaus for three consecutive epochs. Early stopping is applied with patience 7 and a maximum of 40 epochs. Inverse-frequency class weights are applied to the cross-entropy loss as

\begin{equation}
w_c = \frac{1}{N_c + \epsilon}
\label{eq:classweight}
\end{equation}

\subsubsection{Event-to-Usage Mapping}

The six user-facing event types are mapped to weighted combinations of the four usage classes as shown in Table~\ref{tab:mapping}.

\begin{table}[htbp]
\caption{Event-to-Usage Mapping Weights}
\begin{center}
\renewcommand{\arraystretch}{1.15}
\begin{tabular}{|l|l|}
\hline
\textbf{Event} & \textbf{Usage classes (weight)} \\
\hline
Casual     & Casual (1.0), Sports (0.3) \\
Office     & Formal (1.0), Casual (0.2) \\
Wedding    & Ethnic (1.0), Formal (0.7) \\
Party      & Casual (1.0), Ethnic (0.5) \\
Date Night & Formal (0.8), Casual (0.6) \\
Gym        & Sports (1.0) \\
\hline
\end{tabular}
\label{tab:mapping}
\end{center}
\end{table}

The outfit-level event score is the mean of three item-level scores

\begin{equation}
\text{event\_score}(e) = \frac{\displaystyle\sum_{c \in \mathcal{U}(e)} w_{e,c} \cdot p_c}{\displaystyle\sum_{c \in \mathcal{U}(e)} w_{e,c}}
\label{eq:eventscore}
\end{equation}

\subsection{Gender Filtering and Final Ranking}

Before generating combinations, items are filtered so only those matching the selected gender or tagged Unisex are included. Outfits are then ranked in descending order by event score, with cosine compatibility used as a stable tiebreaker

\begin{equation}
\text{rank} = \text{argsort\_desc}(\text{event\_score})
\label{eq:rank}
\end{equation}

% ─────────────────────────────────────────────────────────────
\section{Streamlit Web Application}
% ─────────────────────────────────────────────────────────────

The user interface is a three-page Streamlit application in \texttt{app/app.py}. On the Upload page, users submit a single clothing image along with a name, category, and gender tag. The hybrid extractor processes the image and the embedding, gender, and thumbnail are saved to SQLite. Both models are cached with \texttt{@st.cache\_resource} so they load only once per session.

The My Wardrobe page shows all stored items grouped into Tops, Bottoms, and Shoes. Each item appears as a thumbnail card with its name, gender tag, and a delete button. The Recommend page lets users pick an event type, a gender filter, and a number of outfits between 1 and 10. After clicking Get Recommendations, the system filters by gender, scores all valid combinations, and displays the top-N ranked cards. Each card shows the three item thumbnails alongside a Confidence score for event suitability and a Compatibility score for cosine similarity. Both scores are colour-coded green for 75\% or above, yellow for 55 to 75\%, and red for below 55\%.

% ─────────────────────────────────────────────────────────────
\section{Results and Evaluation}
% ─────────────────────────────────────────────────────────────

\subsection{Event Classifier Performance}

Table~\ref{tab:results} shows per-class precision, recall, and F1 on the held-out test set of 900 items.

\begin{table}[htbp]
\caption{Event Classifier Test Set Results (ResNet50 + CLIP)}
\begin{center}
\renewcommand{\arraystretch}{1.2}
\begin{tabular}{|l|c|c|c|r|}
\hline
\textbf{Class} & \textbf{Prec.} & \textbf{Recall} & \textbf{F1} & \textbf{Support} \\
\hline
Casual  & 0.93 & 0.90 & 0.91 & 615 \\
Ethnic  & 0.86 & 0.89 & 0.88 &  64 \\
Formal  & 0.80 & 0.93 & 0.86 &  75 \\
Sports  & 0.75 & 0.77 & 0.76 & 146 \\
\hline
\textbf{Weighted Avg} & \textbf{0.88} & \textbf{0.88} & \textbf{0.88} & \textbf{900} \\
\hline
\end{tabular}
\label{tab:results}
\end{center}
\end{table}

Overall test accuracy reached 88.0\% with a macro AUC-ROC of 0.9631 and weighted F1 of 0.88. Formal recall of 0.93 and Ethnic recall of 0.89 are particularly important for office and wedding event types and both came in strong. Sports recall is lower at 0.77, most likely because gym-style clothing and casual sportswear are visually quite similar and the model sometimes confuses the two.

\subsection{Backbone Ablation}

Table~\ref{tab:ablation} compares test accuracy across three embedding configurations on the same dataset and training setup.

\begin{table}[htbp]
\caption{Backbone Ablation on Test Accuracy}
\begin{center}
\renewcommand{\arraystretch}{1.2}
\begin{tabular}{|l|c|c|}
\hline
\textbf{Backbone} & \textbf{Dim} & \textbf{Test Acc.} \\
\hline
ResNet50 only          & 2051 & 87.0\% \\
CLIP ViT-B/32 only     &  515 & 87.0\% \\
ResNet50 + CLIP (ours) & 2563 & \textbf{88.0\%} \\
\hline
\end{tabular}
\label{tab:ablation}
\end{center}
\end{table}

The hybrid embedding improves accuracy by 1\% over either backbone alone. While this might seem like a small margin, it confirms the intuition that ResNet texture features and CLIP semantic features are genuinely complementary and that combining them is better than choosing one.

\subsection{Training Dynamics}

The model hit 80\% validation accuracy by epoch 5, which is a fairly quick start. It peaked at 89.3\% around epoch 30 and early stopping kicked in at epoch 37. Validation loss stayed below 0.45 throughout and training loss came down to around 0.01 by epoch 20, which suggests the dropout and weight decay were doing a decent job of keeping overfitting under control.

\subsection{System Performance}

A single image embedding takes roughly 0.2 seconds on GPU running both ResNet50 and CLIP. Scoring 1,000 outfit combinations takes under 0.1 seconds. SQLite reads for 200 items per category are under 10 milliseconds. The embedding migration gate re-embeds 6 wardrobe items in under 3 seconds on GPU.

% ─────────────────────────────────────────────────────────────
\section{Discussion}
% ─────────────────────────────────────────────────────────────

The decision to combine ResNet50 and CLIP was motivated directly by the literature survey findings. ResNet50 alone is a strong visual feature extractor that has been validated across many fashion tasks. But the literature also showed that pure visual similarity is not enough for event-aware recommendations, which is where CLIP's semantic awareness helps. Using both together proved to be better than either alone, as the ablation confirms.

Class merging was preferred over class dropping because dropping classes loses training signal even if those classes are small. Merging Smart Casual into Formal and Party into Casual is a reasonable approximation given that these classes were too small for reliable learning on their own.

Gender is treated as both a feature and a hard filter. As a feature it lets the classifier learn style differences between men's and women's clothing within the same usage class. As a filter it provides the practical guarantee that outfit combinations will be gender-coherent. The filter is the more important of the two from a user perspective.

One thing worth acknowledging is that the event-to-usage mapping in Table~\ref{tab:mapping} is hand-crafted rather than learned from data. Events like Date Night and Party do not have direct equivalents in the training data, so their mappings are an approximation. Better mappings could in principle be learned if occasion-labelled outfit data were available, which comes back to the dataset gap identified in the critical analysis.

The main limitation of the system is that cosine similarity is an unsupervised proxy for compatibility. It works reasonably well in practice because visually similar items tend to coordinate, but it can fail in cases where stylistic contrast is intentional, for example a bright accessory against a neutral outfit. Addressing this would require outfit-level compatibility labels, which do not currently exist in any open dataset.

% ─────────────────────────────────────────────────────────────
\section{Conclusion}
% ─────────────────────────────────────────────────────────────

This paper surveyed the existing literature on fashion recommendation, visual embeddings, and outfit compatibility, and presented Smart Wardrobe as a system that addresses gaps identified in that survey. The key takeaways from the literature were that existing systems lack occasion awareness, do not account for gender, and depend on the Polyvore dataset which is no longer openly accessible. Smart Wardrobe was designed to address all three of these limitations using publicly available data and locally deployable models.

The system achieves 88\% test accuracy and AUC-ROC of 0.9631 on the four-class wardrobe\_v2 event classifier, and the hybrid ResNet50 plus CLIP embedding outperforms either backbone individually. The gender-aware filtering ensures outfit coherence in practice, and the embedding migration gate makes the system robust to future model updates.

Future directions are fairly clear from the limitations identified. A supervised compatibility model trained on labelled outfit pairs would improve on cosine similarity. Personalisation based on user feedback history would make the recommendations more tailored over time. Fine-tuning the CLIP encoder on fashion-specific data could further improve event classification. And extending the item categories beyond tops, bottoms, and shoes to include outerwear and accessories would make the system considerably more practical for real wardrobes.

\section*{Acknowledgment}

This project was built with assistance from Claude (Anthropic) for system architecture, code implementation, and documentation. All AI assistance is explicitly acknowledged here and in the project repository.

\begin{thebibliography}{00}

\bibitem{b1}
K. He, X. Zhang, S. Ren, and J. Sun, ``Deep residual learning for image recognition,'' in \textit{Proc. IEEE Conf. Comput. Vis. Pattern Recognit. (CVPR)}, Las Vegas, NV, Jun. 2016, pp. 770--778.

\bibitem{b2}
X. Han, Z. Wu, Y.-G. Jiang, and L. S. Davis, ``Learning fashion compatibility with bidirectional LSTMs,'' in \textit{Proc. 25th ACM Int. Conf. Multimedia}, Mountain View, CA, Oct. 2017, pp. 1078--1086.

\bibitem{b3}
M. I. Vasileva, B. A. Plummer, K. Dusad, S. Rajpal, R. Kumar, and D. Forsyth, ``Learning type-aware embeddings for fashion compatibility,'' in \textit{Proc. Eur. Conf. Comput. Vis. (ECCV)}, Munich, Sep. 2018, pp. 390--405.

\bibitem{b4}
P. Tangseng, K. Yamaguchi, and T. Okatani, ``Recommending outfits from personal closet,'' in \textit{Proc. IEEE Int. Conf. Comput. Vis. Workshops (ICCVW)}, Venice, Oct. 2017, pp. 2079--2088.

\bibitem{b5}
R. He and J. McAuley, ``VBPR: Visual Bayesian personalized ranking from implicit feedback,'' in \textit{Proc. 30th AAAI Conf. Artif. Intell.}, Phoenix, AZ, Feb. 2016, pp. 144--150.

\bibitem{b6}
P. Aggarwal, ``Fashion product images dataset,'' Kaggle, 2019. [Online]. Available: https://www.kaggle.com/datasets/paramaggarwal/fashion-product-images-dataset

\bibitem{b7}
D. P. Kingma and J. Ba, ``Adam: A method for stochastic optimization,'' in \textit{Proc. 3rd Int. Conf. Learn. Representations (ICLR)}, San Diego, CA, May 2015.

\bibitem{b8}
A. Radford, J. W. Kim, C. Hallacy, A. Ramesh, G. Goh, S. Agarwal, G. Sastry, A. Askell, P. Mishkin, J. Clark, G. Krueger, and I. Sutskever, ``Learning transferable visual models from natural language supervision,'' in \textit{Proc. 38th Int. Conf. Mach. Learn. (ICML)}, Jul. 2021, pp. 8748--8763.

\end{thebibliography}

\end{document}
